\section*{अध्याय \ref{ch:g7:heart-and-circulation} --- दिल और रक्त का चक्कर}

\begin{solution}{exo:g7:heart-and-circulation:1}
दायाँ \omterm{def:g7:heart-and-circulation:rooms}{अलिंद}: शरीर से लौटता \omterm{prop:g4:journey-of-food:absorb}{रक्त} लेता है। दायाँ \omterm{def:g7:heart-and-circulation:rooms}{निलय}: उसे फेफड़ों तक
पंप करता है। बायाँ \omterm{def:g7:heart-and-circulation:rooms}{अलिंद}: फेफड़ों का \omterm{def:g7:what-is-respiration:respiration}{ऑक्सीजन} से भरपूर \omterm{prop:g4:journey-of-food:absorb}{रक्त} लेता है।
बायाँ \omterm{def:g7:heart-and-circulation:rooms}{निलय}: उसे पूरे शरीर भर पंप करता है --- और चारों में सबसे मोटी
दीवार वाला।
\end{solution}

\begin{solution}{exo:g7:heart-and-circulation:2}
वे एकतरफ़ा दरवाज़े हैं, जो \omterm{prop:g4:journey-of-food:absorb}{रक्त} को उलटा उछलने से रोकते हैं। लब-डब उनके
जोड़ी में बंद होने की आवाज़ है: भीतर आने वाले \omterm{def:g7:heart-and-circulation:rooms}{कपाट} तब चटकते हैं जब \omterm{def:g7:heart-and-circulation:rooms}{निलय}
धकेलते हैं (लब), और बाहर जाने वाले \omterm{def:g7:heart-and-circulation:rooms}{कपाट} तब जब वे ढीले पड़ते हैं (डब)।
\end{solution}

\begin{solution}{exo:g7:heart-and-circulation:3}
फेफड़ों वाला चक्कर: दाएँ \omterm{def:g7:heart-and-circulation:rooms}{निलय} से फेफड़ों की \omterm{def:g7:heart-and-circulation:vessels}{केशिकाओं} तक (\omterm{def:g7:what-is-respiration:respiration}{ऑक्सीजन} लादो,
\omterm{def:g7:what-is-respiration:respiration}{कार्बन डाइऑक्साइड} उतारो), और वापस बाएँ \omterm{def:g7:heart-and-circulation:rooms}{अलिंद} तक। शरीर वाला चक्कर: बाएँ
\omterm{def:g7:heart-and-circulation:rooms}{निलय} से पूरे शरीर की \omterm{def:g7:heart-and-circulation:vessels}{केशिकाओं} तक (पहुँचाओ, बटोरो), और वापस दाएँ \omterm{def:g7:heart-and-circulation:rooms}{अलिंद}
तक।
\end{solution}

\begin{solution}{exo:g7:heart-and-circulation:4}
\omterm{def:g7:heart-and-circulation:vessels}{धमनियाँ}: मोटी, लचीली दीवारें --- \omterm{prop:g4:journey-of-food:absorb}{रक्त} को दबाव में \omterm{def:g4:heart-and-blood:heart}{दिल} से दूर ले जाती
हैं। \omterm{def:g7:heart-and-circulation:vessels}{केशिकाएँ}: एक \omterm{def:g5:first-look-at-cells:cell}{कोशिका} भर पतली दीवारें --- यानी हर अदला-बदली की जगह।
\omterm{def:g7:heart-and-circulation:vessels}{शिराएँ}: चौड़ी, पतली दीवारें और एकतरफ़ा \omterm{def:g7:heart-and-circulation:rooms}{कपाट} --- \omterm{prop:g4:journey-of-food:absorb}{रक्त} को कम दबाव पर \omterm{def:g4:heart-and-blood:heart}{दिल}
की ओर लौटाती हैं।
\end{solution}

\begin{solution}{exo:g7:heart-and-circulation:5}
वह शरीर वाला चक्कर चलाता है --- यानी कहीं बड़ा फंदा, पैर की उँगलियों
और \omterm{def:g5:brain-in-command:system}{मस्तिष्क} तक और वापस --- इसलिए उसे सबसे ज़ोर से धकेलना पड़ता है, और
उसकी दीवार उसी हिसाब से बनी है।
\end{solution}

\begin{solution}{exo:g7:heart-and-circulation:6}
\omterm{def:g7:heart-and-circulation:vessels}{केशिकाओं} में। फेफड़ों की \omterm{def:g7:heart-and-circulation:vessels}{केशिकाएँ}: \omterm{def:g7:what-is-respiration:respiration}{ऑक्सीजन} भीतर, \omterm{def:g7:what-is-respiration:respiration}{कार्बन डाइऑक्साइड}
बाहर; \omterm{prop:g7:digestion-and-nutrients:absorption}{रसांकुरों} की: \omterm{def:g7:digestion-and-nutrients:families}{पोषक तत्व} भीतर; \omterm{def:g3:bones-and-muscles:muscle}{पेशियों} की: ईंधन और \omterm{def:g7:what-is-respiration:respiration}{ऑक्सीजन} उतरे,
\omterm{def:g7:what-is-respiration:respiration}{कार्बन डाइऑक्साइड} चढ़ी (और \omterm{def:g1:the-five-senses:sense}{त्वचा} की: ऊष्मा बाहर)।
\end{solution}

\begin{solution}{exo:g7:heart-and-circulation:7}
दायाँ \omterm{def:g7:heart-and-circulation:rooms}{निलय} --- फेफड़ों की \omterm{def:g7:heart-and-circulation:vessels}{केशिकाएँ} (\omterm{def:g7:what-is-respiration:respiration}{ऑक्सीजन} सवार) --- बायाँ \omterm{def:g7:heart-and-circulation:rooms}{अलिंद} ---
बायाँ \omterm{def:g7:heart-and-circulation:rooms}{निलय} --- गरदन की \omterm{def:g7:heart-and-circulation:vessels}{धमनियाँ} --- \omterm{def:g5:brain-in-command:system}{मस्तिष्क} की \omterm{def:g7:heart-and-circulation:vessels}{केशिकाएँ} --- गरदन की
\omterm{def:g7:heart-and-circulation:vessels}{शिराएँ} --- दायाँ \omterm{def:g7:heart-and-circulation:rooms}{अलिंद}। \omterm{def:g5:brain-in-command:system}{मस्तिष्क} की अदला-बदली \omterm{def:g7:what-is-respiration:respiration}{ऑक्सीजन} तथा \omterm{prop:g7:organs-at-work:bill}{ग्लूकोज़} का
एक बड़ा, टिकाऊ हिस्सा लेती है और \omterm{def:g7:what-is-respiration:respiration}{कार्बन डाइऑक्साइड} छोड़ जाती है ---
यानी उस माँग करने वाले सेनापति का पक्का फ़रमान।
\end{solution}

\begin{solution}{exo:g7:heart-and-circulation:8}
\omterm{prop:g4:heart-and-blood:pulse}{नाड़ी} हर \omterm{def:g7:heart-and-circulation:rooms}{निलय} वाली लहर पर \omterm{def:g7:heart-and-circulation:vessels}{धमनी} की दीवार का लचीला तनना है --- यानी ऊँचा
दबाव, और उसके सामने एक तनने वाली मोटी दीवार। \omterm{def:g7:heart-and-circulation:vessels}{शिराओं} तक पहुँचते-पहुँचते
वह लहर \omterm{def:g7:heart-and-circulation:vessels}{केशिकाएँ} पार करने में चुक जाती है: चौड़ी नरम नली में कम, बहा हुआ
दबाव --- धड़कने को कुछ बचता ही नहीं।
\end{solution}

\begin{solution}{exo:g7:heart-and-circulation:9}
वह चढ़ाई वाली वापसी: \omterm{def:g1:my-body:parts}{टाँगों} से \omterm{def:g7:heart-and-circulation:vessels}{शिराओं} का \omterm{prop:g4:journey-of-food:absorb}{रक्त} कम दबाव पर गुरुत्व के
ख़िलाफ़ चढ़ना होता है, और \omterm{def:g7:heart-and-circulation:rooms}{कपाट} हर चढ़ी हुई कमाई थामे रखते हैं। \omterm{def:g1:plants-around-us:parts}{जड़} खड़े
सिपाही की पिंडली की \omterm{def:g3:bones-and-muscles:muscle}{पेशियाँ} \omterm{def:g7:heart-and-circulation:vessels}{शिराओं} को कभी दबाती ही नहीं, चढ़ाई ठहर जाती
है, \omterm{prop:g4:journey-of-food:absorb}{रक्त} \omterm{def:g1:my-body:parts}{टाँगों} में जमा होने लगता है और \omterm{def:g5:brain-in-command:system}{मस्तिष्क} की आपूर्ति गिर जाती है
--- यानी बेहोशी। क़दमताल करते पैर \omterm{def:g7:heart-and-circulation:vessels}{शिराओं} के पंप हैं।
\end{solution}

\begin{solution}{exo:g7:heart-and-circulation:10}
शरीर वाले चक्कर से \omterm{def:g7:heart-and-circulation:vessels}{धमनी} भीतर आती है; भीतर \omterm{prop:g7:digestion-and-nutrients:absorption}{रसांकुरों} की \omterm{def:g7:heart-and-circulation:vessels}{केशिकाएँ} खाने के
\omterm{def:g7:digestion-and-nutrients:families}{पोषक तत्व} उठा लेती हैं
(\cref{prop:g7:digestion-and-nutrients:absorption}); और ख़ास जोड़ यह है
कि उसकी \omterm{def:g7:heart-and-circulation:vessels}{शिरा} पहले छाँटने और जमा करने के लिए \omterm{prop:g7:digestion-and-nutrients:absorption}{यकृत} तक जाती है, और तभी वह
\omterm{prop:g4:journey-of-food:absorb}{रक्त} दाएँ \omterm{def:g7:heart-and-circulation:rooms}{अलिंद} की ओर लौटती आम धारा में मिलता है। नक़्शा \omterm{def:g4:journey-of-food:digestion}{पाचन} का अध्याय
दोबारा सुना देता है।
\end{solution}

\begin{solution}{exo:g7:heart-and-circulation:11}
कमरों का \omterm{prop:g4:journey-of-food:absorb}{रक्त} तो बस बहता हुआ निकल जाता है; और \omterm{def:g4:heart-and-blood:heart}{दिल} की अपनी मोटी दीवारों
को, जो बिना आराम काम करती हैं, किसी भी \omterm{def:g3:bones-and-muscles:muscle}{पेशी} की तरह अपनी \omterm{def:g7:heart-and-circulation:vessels}{केशिका} वाली
सप्लाई चाहिए --- इसीलिए छोटी खिलाने वाली \omterm{def:g7:heart-and-circulation:vessels}{धमनियाँ} ख़ुद \omterm{def:g4:heart-and-blood:heart}{दिल} पर शाखाएँ
बनाती हैं। उन पर परत जम जाए तो वे पंप की ही ईंधन-लाइन का गला घोंट देती
हैं: जो अंग कभी ठहर ही नहीं सकता, उसी को उनकी रुकावट भूखा मारती है।
\end{solution}

\begin{solution}{exo:g7:heart-and-circulation:12}
जिन अदला-बदलियों पर शरीर जीता है --- \omterm{def:g7:lungs-and-gas-exchange:alveoli}{कूपिकाओं} पर \omterm{def:g7:what-is-respiration:respiration}{ऑक्सीजन}, \omterm{prop:g7:digestion-and-nutrients:absorption}{रसांकुरों} पर
\omterm{def:g7:digestion-and-nutrients:families}{पोषक तत्व}, हर \omterm{def:g3:bones-and-muscles:muscle}{पेशी} पर ईंधन और कचरा --- वे सब सिर्फ़ \omterm{def:g7:heart-and-circulation:vessels}{केशिकाओं} की
दीवारों के पार होती हैं; \omterm{def:g7:heart-and-circulation:vessels}{धमनियाँ} और \omterm{def:g7:heart-and-circulation:vessels}{शिराएँ} बस \omterm{prop:g4:journey-of-food:absorb}{रक्त} को उन दीवारों तक लाती
और वहाँ से ले जाती हैं, और \omterm{def:g4:heart-and-blood:heart}{दिल} बस उसे चलता रखता है। पंप और राजमार्ग
उन्हीं बाल जितनी महीन गलियों की ख़ातिर हैं: कहीं भी \omterm{def:g7:heart-and-circulation:vessels}{केशिका} की आवाजाही
रोक दो, और नलसाज़ी चाहे कितनी ही दुरुस्त हो, वह अंग मर जाता है।
\end{solution}

\begin{solution}{pb:g7:heart-and-circulation:1}
\textbf{1.} \omterm{prop:g4:heart-and-blood:pulse}{नाड़ी}: \omterm{def:g7:heart-and-circulation:vessels}{धमनियों} की लहरों की गिनती --- यानी \omterm{def:g1:my-body:joint}{कलाई} पर पढ़ी गई
पंप की दर। रक्तचाप: लहर पर और लहरों के बीच \omterm{def:g7:heart-and-circulation:vessels}{धमनी} का दबाव --- यानी
राजमार्गों पर पड़ा बोझ। स्टेथस्कोप की लब-डब: \omterm{def:g7:heart-and-circulation:rooms}{कपाट} बंद होते हुए सुने गए
--- यानी दरवाज़ों की दुरुस्ती।

\textbf{2.} चुस्त मरीज़ में कम आराम वाली \omterm{prop:g4:heart-and-blood:pulse}{नाड़ी} का मतलब है एक मज़बूत
\omterm{def:g7:heart-and-circulation:rooms}{निलय}, जो हर धड़कन पर ज़्यादा \omterm{prop:g4:journey-of-food:absorb}{रक्त} चलाता है: प्रशिक्षित \omterm{def:g4:heart-and-blood:heart}{दिल} शांत सुस्ताता
है, और उसके पास बचत भी होती है।

\textbf{3.} \omterm{def:g7:heart-and-circulation:rooms}{कपाटों} के दरवाज़े: \omterm{def:g7:heart-and-circulation:rooms}{निलयों} के धकेलने पर भीतर आने वाले \omterm{def:g7:heart-and-circulation:rooms}{कपाट}
(लब), और उनके ढीले पड़ने पर बाहर जाने वाले \omterm{def:g7:heart-and-circulation:rooms}{कपाट} (डब) --- यानी चार
दरवाज़े, दो जोड़ियों में, साफ़-साफ़ बंद होते हुए।

\textbf{4.} दायाँ \omterm{def:g7:heart-and-circulation:rooms}{अलिंद}, दायाँ \omterm{def:g7:heart-and-circulation:rooms}{निलय}, फेफड़ों की \omterm{def:g7:heart-and-circulation:vessels}{केशिकाएँ}, बायाँ \omterm{def:g7:heart-and-circulation:rooms}{अलिंद},
बायाँ \omterm{def:g7:heart-and-circulation:rooms}{निलय}, शरीर की \omterm{def:g7:heart-and-circulation:vessels}{धमनियाँ}, शरीर की \omterm{def:g7:heart-and-circulation:vessels}{केशिकाएँ} (मान लो, किसी \omterm{def:g3:bones-and-muscles:muscle}{पेशी} की),
\omterm{def:g7:heart-and-circulation:vessels}{शिराएँ} --- और घर वापस दाएँ \omterm{def:g7:heart-and-circulation:rooms}{अलिंद} में।

\textbf{5.} हर धड़कन अपनी लहर का एक हिस्सा उस रिसाव से उलटा गँवा देती
है, इसलिए शरीर वाले चक्कर का हर धड़कन पर पहुँचाना घट जाता है; \omterm{def:g4:heart-and-blood:heart}{दिल}
ज़्यादा धड़कनों और ज़्यादा मेहनत से भरपाई करता है, और दमख़म गिर जाता है
--- वह घरघराहट उलटी बही हुई डिलीवरी है।

\textbf{6.} रक्तचाप चढ़ा हुआ (अकड़े, सिकुड़े राजमार्ग लहर का विरोध करते
हैं) और \omterm{prop:g4:heart-and-blood:pulse}{नाड़ी} वाला \omterm{def:g7:heart-and-circulation:vessels}{धमनी} का तनना बदला हुआ; और बायाँ \omterm{def:g7:heart-and-circulation:rooms}{निलय}, जो हर चक्कर पर
अकड़े हुए परिपथ के ख़िलाफ़ धकेलता है, ओवरटाइम करता और मोटा होता जाता
है।

\textbf{7.} \omterm{def:g4:heart-and-blood:heart}{दिल} की अपनी खिलाने वाली \omterm{def:g7:heart-and-circulation:vessels}{धमनियों} में परत जमना उस \omterm{def:g3:bones-and-muscles:muscle}{पेशी} की
ईंधन-लाइन का गला घोंटता है जो कभी आराम कर ही नहीं सकती --- और यही \omterm{def:g4:heart-and-blood:heart}{दिल}
के दौरे तक का जाना-पहचाना रास्ता है।

\textbf{8.} \omterm{def:g7:heart-and-circulation:vessels}{शिराओं} के एकतरफ़ा \omterm{def:g7:heart-and-circulation:rooms}{कपाट}, और हरकत वाला वह पेशी-पंप जो ग़ायब
है: \omterm{def:g1:plants-around-us:parts}{जड़} खड़े रहने पर पिंडलियाँ कभी दबाती ही नहीं, \omterm{def:g7:heart-and-circulation:rooms}{कपाटों} की कमाई ठहर
जाती है, और टख़नों में जमाव हो जाता है। चलना \omterm{def:g7:heart-and-circulation:vessels}{शिराओं} को लय में दबाता है
--- यानी वह राहत यांत्रिक है।

\textbf{9.} क्योंकि शरीर का दौरा \omterm{prop:g4:journey-of-food:absorb}{रक्त} की \omterm{def:g7:what-is-respiration:respiration}{ऑक्सीजन} ख़र्च कर देता है:
उसे सिर्फ़ फेफड़ों वाला चक्कर ही दोबारा लाद सकता है। \omterm{def:g7:how-animals-breathe:four}{फेफड़े} छोड़ देने का
मतलब होता चुके हुए \omterm{prop:g4:journey-of-food:absorb}{रक्त} को ख़ाली गाड़ी के साथ पहुँचाने भेजना --- दोहरा
फंदा ``पहले भरो, फिर बाँटो'' है, और उसे \omterm{def:g4:heart-and-blood:heart}{दिल} के दोनों पहलू लागू करते
हैं।

\textbf{10.} \omterm{def:g7:heart-and-circulation:vessels}{धमनी}: मोटी, लचीली, ऊँचा दबाव --- यानी बाहर जाते राजमार्ग।
\omterm{def:g7:heart-and-circulation:vessels}{केशिका}: एक \omterm{def:g5:first-look-at-cells:cell}{कोशिका} भर पतली, हर चीज़ में पिरोई हुई --- यानी पहुँचाने वाली
गलियाँ, जहाँ सारी अदला-बदलियाँ होती हैं। \omterm{def:g7:heart-and-circulation:vessels}{शिरा}: चौड़ी, \omterm{def:g7:heart-and-circulation:rooms}{कपाटों} वाली, कम
दबाव --- यानी वापसी।

\textbf{11.} हर परामर्श की तकलीफ़ ठीक वहीं मायने रखी जहाँ उसने \omterm{def:g7:heart-and-circulation:vessels}{केशिका}
की सेवा को छुआ: रिसाव ने पहुँचाना भूखा रखा, परत ने आपूर्ति की लाइनें गला
घोंटीं, और जमाव ने वापसी ठहरा दी। पढ़तें, \omterm{def:g7:heart-and-circulation:rooms}{कपाट} और राजमार्ग साधन हैं;
\omterm{def:g7:heart-and-circulation:vessels}{केशिका} की आवाजाही --- यानी अदला-बदलियाँ --- ही मंज़िल है।

\textbf{12.} रोज़ चलो-फिरो --- पंप मज़बूत होता है और राजमार्ग लचीले
रहते हैं; संतुलित खाओ --- \omterm{def:g7:heart-and-circulation:vessels}{धमनियों} की दीवारें परत से साफ़ रहती हैं; और
धुआँ बिलकुल नहीं --- \omterm{def:g4:heart-and-blood:vessels}{वाहिकाएँ} न सिकुड़ती हैं न अकड़ती हैं, और \omterm{prop:g4:journey-of-food:absorb}{रक्त} की
सीटें \omterm{def:g7:what-is-respiration:respiration}{ऑक्सीजन} के लिए ख़ाली रहती हैं।
\end{solution}
